\documentclass[10pt,aspectratio=169]{beamer}
\usepackage[utf8]{inputenc}
\usepackage{amsmath,amssymb,amsfonts,bm,mathtools}
\usepackage{booktabs}

% Personalización del Tema Beamer
\usetheme{Madrid}

% Paleta de Colores Institucional IPN
% Guinda oficial: Pantone 222 C, RGB #750946
% Gris oficial: Pantone Cool Gray 10C, RGB #636569
\definecolor{IPNGuinda}{RGB}{117, 9, 70}      % #750946
\definecolor{IPNGris}{RGB}{99, 101, 105}       % #636569
\definecolor{IPNNegro}{RGB}{35, 31, 32}        % #231f20
\definecolor{LightBg}{RGB}{245, 247, 250}

\setbeamercolor{palette primary}{bg=IPNGuinda,fg=white}
\setbeamercolor{palette secondary}{bg=IPNGris,fg=white}
\setbeamercolor{palette tertiary}{bg=IPNGuinda!90,fg=white}
\setbeamercolor{titlelike}{parent=palette primary}
\setbeamercolor{block title}{bg=IPNGuinda,fg=white}
\setbeamercolor{block body}{bg=LightBg,fg=black}
\setbeamercolor{block title example}{bg=IPNGris,fg=white}
\setbeamercolor{block body example}{bg=LightBg,fg=black}
\setbeamercolor{block title alert}{bg=IPNNegro,fg=white}
\setbeamercolor{block body alert}{bg=LightBg,fg=black}

\title[Control del Péndulo de Furuta]{Modelado y Control del Péndulo de Furuta}
\subtitle{Un enfoque práctico con espacios funcionales y control embebido}
\author[Propuesta Doctoral]{Propuesta de Avance - Investigación Doctoral}
\institute[Control y Matemáticas Aplicadas]{Reunión 1 - Comité tutorial}
\date{\today}

\begin{document}

% --- SLIDE 1 ---
\begin{frame}
  \titlepage
  \begin{center}
    \small \textit{"Fundamentado en el rigor analítico de la teoría de control clásico y el análisis funcional."}
  \end{center}
\end{frame}

% --- SLIDE 2 ---
\begin{frame}{2. Marco Metodológico de Referencia: Flujo de 10 Etapas}
  \begin{block}{Metodología Base de Síntesis de Control e IA}
    Se adopta el marco metodológico (Z.P. et. al.) de 10 etapas para diseño de controladores en tiempo real.
 Se cuenta con un prototipo físico inicial (plataforma Furuta instrumentada con microcontrolador, encoders, etc), lo que permite que la propuesta se centre en la validación formal e implementación de controlador(es),  bajo una ruta de trabajo ordenada y verificable.
  \end{block}
  
  \vspace{0.2cm}
  \begin{columns}[T]
    \begin{column}{0.48\textwidth}
      \textbf{Etapas Fundamentales de la Metodología:}
      \begin{enumerate}
        \setcounter{enumi}{0}
        \item \small \textbf{Planteamiento del problema} y especificaciones físicas.
        \item \small \textbf{Modelado del sistema} (Euler-Lagrange y State-Space).
        \item \small \textbf{Control aproximado / IA} (Inferencia difusa ANFIS).
        \item \small \textbf{Simplificación funcional} (Aproximación polinomial).
        \item \small \textbf{Generación de Código C} estático.
      \end{enumerate}
    \end{column}
    \begin{column}{0.48\textwidth}
      \begin{enumerate}
        \setcounter{enumi}{5}
        \item \small \textbf{Implementación embebida} en microcontrolador.
        \item \small \textbf{Pruebas en Hardware-in-the-Loop (HIL)}.
        \item \small \textbf{Sintonía fina y métricas} ($\|e\|_2, \|e\|_\infty$).
        \item \small \textbf{Validación experimental} en planta real.
        \item \small \textbf{Documentación y transferencia}.
      \end{enumerate}
    \end{column}
  \end{columns}
\end{frame}

% --- SLIDE 3 ---
\begin{frame}{3. Dinámica Cinemática No Lineal (Euler-Lagrange)}
  \begin{block}{Modelo Físico del Péndulo de Furuta (2 GDL Subactuado)}
    Sea $q = [\theta_1, \theta_2]^T$ el vector de coordenadas generalizadas (brazo rotatorio $\theta_1$, péndulo $\theta_2$). La ecuación vectorial de movimiento está dada por:
    \begin{equation*}
      \mathbf{M}(q)\ddot{q} + \mathbf{C}(q,\dot{q})\dot{q} + \mathbf{G}(q) + \mathbf{F}_{fric}(\dot{q}) = \mathbf{B} u
    \end{equation*}
  \end{block}
  
  \vspace{0.1cm}
  \begin{itemize}
    \item \textbf{Matriz de Inercia Simétrica y Definida Positiva:} 
    $\mathbf{M}(q) = \begin{bmatrix} J_1 + m_2 L_1^2 + m_2 r_2^2 \sin^2\theta_2 & -m_2 L_1 r_2 \cos\theta_2 \\ -m_2 L_1 r_2 \cos\theta_2 & J_2 + m_2 r_2^2 \end{bmatrix}$
    \item \textbf{Términos de Coriolis y Centrífugos:} Acoplamiento no lineal altamente persistente $\mathbf{C}(q, \dot{q})$.
    \item \textbf{Subactuación Extrema:} $\mathbf{B} = [1, 0]^T$; sólo el brazo horizontal recibe par directo $u$.
  \end{itemize}
\end{frame}

% --- SLIDE 4 ---
\begin{frame}{4. Linealización y Análisis en Espacio de Estados}
  \begin{block}{Representación Matricial en el Punto de Equilibrio Inestable}
    Linealizando en el equilibrio inestable $x^* = [0, \pi, 0, 0]^T$ con $x = [\theta_1, \theta_2, \dot{\theta}_1, \dot{\theta}_2]^T$:
    \begin{equation*}
      \dot{x}(t) = \mathbf{A} x(t) + \mathbf{B} u(t), \quad y(t) = \mathbf{C} x(t)
    \end{equation*}
  \end{block}

  \begin{columns}[T]
    \begin{column}{0.48\textwidth}
      \begin{exampleblock}{Controlabilidad ($\mathcal{C}$)}
        \begin{equation*}
          \mathcal{C} = \begin{bmatrix} \mathbf{B} & \mathbf{A}\mathbf{B} & \mathbf{A}^2\mathbf{B} & \mathbf{A}^3\mathbf{B} \end{bmatrix}
        \end{equation*}
        $\text{rank}(\mathcal{C}) = 4 \implies$ \textbf{Modos inestables completamente alcanzables.}
      \end{exampleblock}
    \end{column}
    \begin{column}{0.48\textwidth}
      \begin{exampleblock}{Observabilidad ($\mathcal{O}$)}
        \begin{equation*}
          \mathcal{O} = \begin{bmatrix} \mathbf{C} \\ \mathbf{C}\mathbf{A} \\ \mathbf{C}\mathbf{A}^2 \\ \mathbf{C}\mathbf{A}^3 \end{bmatrix}
        \end{equation*}
        $\text{rank}(\mathcal{O}) = 4 \implies$ \textbf{Reconstrucción de estados por observador de Luenberger.}
      \end{exampleblock}
    \end{column}
  \end{columns}
\end{frame}

% --- SLIDE 5 ---
\begin{frame}{5. Fenómenos Reales No Modelados: Fricción y Retardo Digital}
  \begin{alertblock}{La Grieta del Modelo Lineal Determinista}
    En la práctica experimental existen dos perturbaciones severas que degradan la estabilidad de reguladores LQR/PID clásicos:
  \end{alertblock}

  \vspace{0.2cm}
  \begin{columns}[T]
    \begin{column}{0.48\textwidth}
      \textbf{1. Fricción No Lineal en Rodamientos:}
      \begin{equation*}
        F_{fric}(\dot{q}) = f_v \dot{q} + f_c \text{sgn}(\dot{q}) + f_s e^{-(\dot{q}/\dot{q}_s)^2}\text{sgn}(\dot{q})
      \end{equation*}
      \small Fenómeno de Stribeck y zona muerta no diferenciable en $\dot{q}=0$.
    \end{column}
    \begin{column}{0.48\textwidth}
      \textbf{2. Retardo Digital de Muestreo ($\tau$):}
      \begin{equation*}
        u(t) = u_{calc}(t - \tau), \quad H(s) = e^{-s\tau}
      \end{equation*}
      \small Desfasaje introducido por adquisición ADC/Encoder y cómputo del microcontrolador.
    \end{column}
  \end{columns}
\end{frame}

% --- SLIDE 6 ---
\begin{frame}{6. Justificación Rigurosa del Control Inteligente}
  \begin{block}{Control Difuso como Aproximador Universal}
    La Inteligencia Artificial no reemplaza la teoría de control; actúa como una cartografía suave sobre la variedad no lineal no estructurada:
    \begin{equation*}
      u_{IA}(x) = \frac{\sum_{i=1}^M w_i(x) f_i(x)}{\sum_{i=1}^M w_i(x)}, \quad w_i(x) = \prod_{j=1}^n \mu_{A_{ji}}(x_j)
    \end{equation*}
  \end{block}

  \begin{itemize}
    \item \textbf{Teorema de Aproximación Universal (Stone-Weierstrass):} Los sistemas de inferencia Sugeno/Mamdani son densos en el espacio de funciones continuas $C(K)$ sobre un compacto $K \subset \mathbb{R}^n$.
    \item \textbf{Propósito en la Propuesta:} Interceptar la incertidumbre paramétrica $\Delta \mathbf{M}(q)$ y compensar la fricción estática $f_c \text{sgn}(\dot{q})$ sin chattering ni discontinuidades.
  \end{itemize}
\end{frame}

% --- SLIDE 7 ---
\begin{frame}{7. Fundamentación en Análisis Funcional (Normas y Operadores)}
  \begin{block}{Certificados Analíticos en Espacios de Banach e Hilbert}
    Consideramos la trayectoria de error $e(t) = x(t) - x^*(t)$ como un elemento del espacio funcional $L_2([0, T])$ o $L_\infty([0, T])$.
  \end{block}

  \vspace{0.2cm}
  \begin{columns}[T]
    \begin{column}{0.48\textwidth}
      \textbf{Norma de Energía del Error ($L_2$):}
      \begin{equation*}
        \|e\|_2 = \left( \int_0^T \|e(t)\|^2 dt \right)^{1/2} < \infty
      \end{equation*}
      \small Garantiza que la energía acumulada del error de balanceo sea finita (Estabilidad en el sentido de Lyapunov).
    \end{column}
    \begin{column}{0.48\textwidth}
      \textbf{Norma de Máxima Desviación ($L_\infty$):}
      \begin{equation*}
        \|e\|_\infty = \text{ess\,sup}_{t \in [0,T]} \|e(t)\| \le \epsilon
      \end{equation*}
      \small Acota la excursión máxima del péndulo para evitar la caída durante perturbaciones impulsivas.
    \end{column}
  \end{columns}
\end{frame}

% --- SLIDE 8 ---
\begin{frame}{8. Arquitectura Híbrida Propuesta}
  \begin{block}{Estrategia de Dos Fases: Swing-Up Eficiente y Balanceo Robusto}
    Integración de control por energía en grandes desviaciones y control óptimo-difuso en la vecindad del equilibrio inestable.
  \end{block}

  \vspace{0.2cm}
  \begin{columns}[T]
    \begin{column}{0.48\textwidth}
      \begin{exampleblock}{Fase 1: Swing-Up por Energía}
        \begin{equation*}
          E(q, \dot{q}) = \frac{1}{2} J_2 \dot{\theta}_2^2 + m_2 g r_2 (1 - \cos\theta_2)
        \end{equation*}
        $u_{swing} = k_E (E - E_0) \text{sgn}(\dot{\theta}_2 \cos\theta_1)$
      \end{exampleblock}
    \end{column}
    \begin{column}{0.48\textwidth}
      \begin{exampleblock}{Fase 2: Balanceo Adaptativo Híbrido}
        \begin{equation*}
          u(t) = -\mathbf{K}_{LQR} \hat{x}(t) + u_{difuso}(\hat{x}, t - \tau)
        \end{equation*}
        Compensación en tiempo real del retardo $\tau$ e incertidumbres de masa/fricción.
      \end{exampleblock}
    \end{column}
  \end{columns}
\end{frame}

% --- SLIDE 9 ---
\begin{frame}{9. Simplificación Funcional y Generación de Código C}
  \begin{block}{Traducción de Expresiones Complejas a Ejecutables de Tiempo Real}
    Para cumplir con los estrictos períodos de muestreo en microcontroladores ($T_s \le 1\text{ ms}$), se simplifica el operador de IA (Etapas 4 y 5 de la Metodología).
  \end{block}

  \begin{itemize}
    \item \textbf{Aproximación Polinomial / Tablas de Búsqueda (LUT):} Proyección del mapa difuso $u_{difuso}(x)$ sobre una superficie polinomial suave de orden reducido o grilla interpolada linealmente.
    \item \textbf{Generación Automática de Código C Determinista:}
      \begin{itemize}
        \item Sin asignación dinámica de memoria (\texttt{malloc} prohibido).
        \item Estructuras fijas en punto flotante (\texttt{float32\_t}) optimizadas para FPU.
        \item Ejecución garantizada en $O(1)$ pasos de cómputo por ciclo de reloj.
      \end{itemize}
  \end{itemize}
\end{frame}

% --- SLIDE 10 ---
\begin{frame}{10. Banco de Pruebas y Arquitectura de Hardware}
  \begin{block}{Infraestructura Experimental del Péndulo de Furuta}
    Validación física del controlador híbrido con instrumentación de alta resolución.
  \end{block}

  \begin{columns}[T]
    \begin{column}{0.48\textwidth}
      \textbf{Sensores y Adquisición:}
      \begin{itemize}
        \item Sensor Magnético Angular \textbf{AS5600} (12 bits, I2C/PWM) en el péndulo vertical.
        \item Encoder óptico incremental \textbf{LPD3806} (600 PPR) en el brazo horizontal.
      \end{itemize}
    \end{column}
    \begin{column}{0.48\textwidth}
      \textbf{Actuación y Cómputo:}
      \begin{itemize}
        \item Etapa de potencia: Puente H \textbf{BTS7960} (43A) con modulación PWM a 20 kHz.
        \item Procesador Embedded: Microcontrolador 32-bit (STM32/ESP32) ejecutando RTOS.
      \end{itemize}
    \end{column}
  \end{columns}
\end{frame}

% --- SLIDE 11 ---
\begin{frame}{11. Cronograma de Investigación Doctoral (6 Semestres)}
  \begin{table}
    \centering
    \small
    \begin{tabular}{lp{8.5cm}}
      \toprule
      \textbf{Periodo} & \textbf{Hitos y Entregables Clave} \\
      \midrule
      Semestres 3--4 & Modelado dinámico riguroso, caracterización de fricción y simulación en python. \\
      Semestres 5--6 &  Estimación del retardo y algoritmos difusos. Simplificación polinomial, traducción a C y pruebas HIL. \\
      Semestres 7--8 & Pruebas experimentales finales en planta física, redacción de tesis y artículos. \\
      \bottomrule
    \end{tabular}
  \end{table}
\end{frame}

% --- SLIDE 12 ---
\begin{frame}{12. Conclusiones y Compromiso de Rigor Académico}
  \begin{block}{Síntesis del Aporte Propuesto}
    Un esquema integral que conecta la abstracción del análisis funcional con la exigencia práctica del control embebido en tiempo real sobre el péndulo de Furuta.
  \end{block}

  \vspace{0.2cm}
  \begin{itemize}
    \item \textbf{Generacion de algos:} La IA no sustituye la física; la enriquece capturando la no linealidad residual.
    \item \textbf{Validez Matemática:} Certificación de estabilidad mediante acotación en normas $L_2$ y $L_\infty$.
    \item \textbf{Factibilidad Técnica:} Demostración de viabilidad en código C estático con hardware accesible.
  \end{itemize}

  \vspace{0.4cm}
  \begin{center}
    \Large \textbf{Muchas gracias por su atención.} \\
    \small \textit{Quedo a su entera disposición para preguntas y comentarios.}
  \end{center}
\end{frame}

\end{document}